\section{Key Claims and Results}

\subsection*{Problem formulation}
The two-layer model is
\begin{equation}
\hat y(x;\theta)=\frac{1}{N}\sum_{i=1}^N \sigma_*(x;\theta_i),
\end{equation}
and the population square-loss risk can be rewritten as
\begin{equation}
R_N(\theta)
=
R_\# + \frac{2}{N}\sum_{i=1}^N V(\theta_i)
+ \frac{2}{N^2}\sum_{i,j=1}^N U(\theta_i,\theta_j).
\end{equation}
This motivates the mean-field functional
\begin{equation}
R(\rho)=R_\# + 2\int V(\theta)\,\rho(\dd\theta)
+ \int U(\theta_1,\theta_2)\,\rho(\dd\theta_1)\rho(\dd\theta_2).
\end{equation}

\subsection*{Main theorem-level results from the paper}

\paragraph{Lemma 1.}
For the isotropic Gaussian example, this gives a criterion for when the uniform measure on a sphere of radius $r_*$ is a global minimizer of the reduced risk. This is the first concrete bridge from the abstract mean-field risk to an explicit optimizer in a model.

\paragraph{Theorem 1.}
In the isotropic Gaussian classification model, for good initialization and large enough dimension / width, SGD reaches risk within $\eta$ of the global optimum after a finite scaled time window. The key message is not just convergence, but convergence with sample complexity linear in dimension and independent of width once width is large enough.

\paragraph{Theorem 2.}
An analogous near-optimal SGD result holds for anisotropic Gaussian classification, where the task involves recovering a relevant subspace. This shows the framework is not restricted to a fully symmetric toy problem.

\paragraph{Theorem 3.}
This is the paper's central structural theorem: SGD and noisy SGD converge to the corresponding PDEs on finite time intervals. Formally, the empirical measure of parameters converges weakly to the solution of DD / diffusion DD. This is the propagation-of-chaos step that converts the optimization analysis into PDE analysis.

\paragraph{Proposition 1.}
The finite-width optimum differs from the mean-field optimum by at most $O(1/N)$ under a mild moment condition. This justifies using $\inf_\rho R(\rho)$ as a proxy for the large-width optimum.

\paragraph{Propositions 2 and 3.}
These characterize fixed points of the noiseless and noisy distributional dynamics. Proposition~2 gives the zero-current condition for noiseless DD. Proposition~3 gives the Boltzmann fixed-point formula
\begin{equation}
\rho_*(\theta)=\frac{1}{Z(\beta)}\exp\{-\beta \Psi_\lambda(\theta;\rho_*)\},
\end{equation}
for the noisy regularized flow.

\paragraph{Theorem 4.}
Under A1--A4 and $\lambda>0$, the free energy has a unique minimizer $\rho_{\beta,\lambda}^*$ and the diffusion PDE converges to it. This is the cleanest generic asymptotic result in the paper.

\paragraph{Theorem 5.}
Regularized noisy SGD inherits the PDE convergence and achieves near-optimal regularized risk. This is the main generic algorithmic corollary.

\paragraph{Theorems 6 and 7.}
These analyze noiseless fixed points. Theorem~6 gives a Hessian-based local stability condition for an atomic fixed point $\delta_{\theta^*}$: positive definite local Hessian implies local exponential convergence. Theorem~7 gives a converse instability-type statement for certain mixed fixed points with a saddle direction.

\subsection*{Free-energy identity}
One of the most important formulas in the paper is the free-energy dissipation law for the noisy PDE:
\begin{equation}
\label{eq:free-energy-note}
\frac{\dd}{\dd t}F_{\beta,\lambda}(\rho_t)
=
-2\xi(t)\int
\bigl\|\nabla_\theta \bigl(\Psi_\lambda(\theta;\rho_t)+\beta^{-1}\log \rho_t(\theta)\bigr)\bigr\|_2^2\,
\rho_t(\theta)\,\dd\theta.
\end{equation}
This is the backbone of Theorem~4. It is the point where the Wasserstein gradient-flow structure becomes operational.
